\section{Introduction}

Phase retrieval aims to recover a complex-valued object from diffraction
measurements whose phase information is not directly observed. In coherent
diffraction imaging, this inverse problem is highly non-linear and admits
ambiguities unless additional physical priors are imposed.

Deep neural networks offer a data-driven alternative to iterative phase
retrieval algorithms. AutoPhaseNN demonstrates that a convolutional neural
network can be trained to map diffraction amplitudes to object-domain amplitude
and phase, followed by a differentiable Fourier-domain reconstruction step.
In this work, AutoPhaseNN is used first as a reproduction target rather than
only as a baseline. This choice is important: before testing a new architecture,
the data preprocessing, Fourier projection, support masking, checkpoint loading,
and evaluation metrics must match the reference pipeline closely enough that
model differences are not confused with implementation differences.

The current experimental evidence suggests a central problem: optimizing
diffraction-domain L1 loss alone does not guarantee a physically meaningful
real-space solution. This motivates a careful comparison between Fourier-domain
metrics and real-space metrics such as amplitude error, phase error, and support
agreement.

The main question is therefore not simply whether a UMamba-based network can
reduce far-field loss. The more useful question is whether it can do so while
recovering an object-domain representation that is consistent with the physical
support and phase structure expected by the phase retrieval problem. This paper
reports the current reproduction results, the UMamba replacement experiments,
and the diagnostic evidence behind this distinction.

The contributions of this draft are:
\begin{itemize}
  \item We establish a reproducible AutoPhaseNN-style PyTorch training and
  evaluation pipeline.
  \item We compare the original AutoPhaseNN backbone and UMamba-based variants
  under aligned data loading, post-processing, checkpointing, and metrics.
  \item We show how far-field losses can be inconsistent with real-space
  reconstruction quality, especially when support structure is weakly enforced.
  \item We organize the evidence for spatial support priors as a key factor in
  stable neural phase retrieval.
\end{itemize}
